\section{Additional pre-specified tests}\label{sec:plan4}

The tests above answer a conditional question and concern one model. A second plan (\repoid{docs/REVISION\_PLAN\_4.md}) was written after the results above were known and before any of the following had been computed, with two exceptions it names: the winsorized mean-loss model confidence set and the firm-block interval of the shortfall test, which the first plan's code had already produced and which Sections~\ref{sec:prespec} and~\ref{sec:pillarA} report. It specifies an unconditional evaluation, a second model under the same protocol and the statistics the first plan omitted. It was committed, timestamped with OpenTimestamps, now anchored in the Bitcoin blockchain, and published before its code was written. Three dated addenda, each timestamped in the same way, record the implementation choices, one bug fix made after the first results had been read, and statistics that the plan lists but the first run had not computed.\footnote{The bug: the density rival and the combination weight of the unconditional test had been trained on test years only, which also reduced the test years to 2020--2024. Before the fix the mean CRPS was 21.7 for the simulator against 15.7 for the rival, and the combination gain was 3.6\% in winsorized mean absolute error. The corrected run is reported.} The simulations ran under the code, data and walk-forward choices of the full run, with one module added. The simulator's baseline health score reproduces that run exactly on all \PRepN{} firm-years. Four tests are confirmatory (Table~\ref{tab:rev4}) and the rest are descriptive. The supplementary material reports every statistic of the plan that this section does not.

\begin{table}[t]
\centering\scriptsize
\caption{Confirmatory tests of the second plan}
\label{tab:rev4}
\resizebox{\linewidth}{!}{\begin{tabular}{llll}
\toprule
Test & Comparison & Statistic [95\% CI] & Verdict \\
\midrule
P4-1 A & Unconditional point: simulator vs margin $\times$ seasonal random walk & 0.95 [0.09, 1.81] & fails \\
P4-1 B & Unconditional density: CRPS simulator vs empirical rival & 19.2 vs 14.1 ($p$ 0.997) & fails \\
P4-2 & Second model (pooled LAD) vs 16-quarter margin, median error & 0.18 [$-$0.50, 0.86] & fails \\
P4-7 & Recalibrated shortfall probability vs logit, label L2, BSS & $-$0.015 [$-$0.026, $-$0.005] & fails \\
\bottomrule
\end{tabular}
}
\par\smallskip
\begin{minipage}{\linewidth}\footnotesize
Point tests: the mean over test years of the yearly difference in median absolute error of eight-quarter operating cash flow (percentage points of trailing revenue), with its $t(G-1)$ interval; success requires the upper bound below zero. Density: mean CRPS, Diebold--Mariano test under firm$\times$year clustering, one-sided. Shortfall: Brier skill score of the isotonically recalibrated simulator against the logit, with a firm-block interval; label L2 excludes delisting notices.
\end{minipage}
\end{table}

\paragraph{Unconditional forecasts} Without the realized revenue, the simulator's own baseline paths give, on the \PUncCommonN{} firm-years of 2017--2024 that all four unconditional forecasts cover, a median absolute error of \PUncCMedTwin{} points. The 16-quarter margin applied to a seasonal random walk of revenue gives \PUncCMedRw, and the four-quarter margin \PUncCMedFour. The simulator is worse in \PUncWorse{} of \PUncG{} years. The mean of the yearly differences is \PUncStatMean{} points, with an interval of \PUncLo{} to \PUncHi, so the test \PUncVerdict. The Diebold--Mariano test gives $p=\PUncDmP$, and the wild-cluster bootstrap by year $p=\PUncWildYearP$ for the yearly differences and $p=\PUncWildP$ for the firm-level ones. The exact randomization test over years gives $p$ \PUncRandPrel{}, and the sign test, which needs eight out of eight years at this level, $p=\PUncSignP$.

The 16-quarter margin applied to the simulator's own median revenue path gives \PUncCMedOwn. That is better than the simulator in every year (interval of the difference \PUncOwnLo{} to \PUncOwnHi). Against the random walk of revenue the interval is \PUncOwnRwLo{} to \PUncOwnRwHi. Under median loss, then, the simulated revenue path is no worse than a random walk, and the simulator's loss arises in its mapping from revenue to cash. Under mean loss the same margin is worse than the random walk (\PUncCMeanOwn{} against \PUncCMeanRw).

The unconditional density test also \PUncDenVerdict. The mean CRPS is \PUncCrpsTwin{} against \PUncCrpsRival{} for the margin with its empirical residuals, and the simulator is better in \PUncDenBetter{} of \PUncDenYears{} years. Its 90\% intervals cover \PUncCovTwin{} of outcomes, against \PUncCovRival{} for the rival. Adding a fraction of \PUncW{} of the simulator's departure to the rival lowers the winsorized mean absolute error by \PUncGainW{} and the median by \PUncGainMed, and changes the raw mean by \PUncGainRaw.

\paragraph{A second model} The second model is a pooled, walk-forward median regression, ML1. It uses the inputs of the margin rules (the three margins and their products with realized revenue), the opening ratios of the shortfall logit, revenue volatility, log revenue, the fixed-cost share from the simulator's calibration and industry. It is trained on cutoffs up to $Y-3$, and it needs 500 training firm-years, so it is tested on 2018--2024. Its median error is \PMlMed{} against \PMlMedNthree{} for the 16-quarter margin. The mean yearly difference is \PMlStatMean{}, and its interval, \PMlLo{} to \PMlHi, includes zero, so the protocol does not declare this model better than persistence either.

Descriptively, ML1 has a significantly smaller median error than the simulator (\PMlTwinMed{} on the same firm-years; interval of the difference \PMlTwinLo{} to \PMlTwinHi) and a smaller mean absolute error than the margin (\PMlMean{} against \PMlMeanNthree). Its predictive quantiles cover \PMlCov{} of outcomes.

In an encompassing regression of the outcome on all three forecasts, winsorized at the 1st and 99th percentiles, the simulator keeps a significant weight ($t=\PMlEncTwinT$; the regression $t=\PMlEncMlT$, the margin $t=\PMlEncNthreeT$; critical value \PMlEncCrit). Its departures therefore carry information about mean cash flow that neither the margin nor the regression has.

In the model confidence set with ML1 added, ML1 and the margin both remain ($p_{\mathrm{MCS}}=\PMcsMlMedMlP$ and \PMcsMlMedNthreeP{} under median loss), and the margin plus the simulator is again the last model standing. The stand-alone simulator is excluded, together with its loop-free version, both cash-revenue regressions and the four-quarter and shrunk margins: $p_{\mathrm{MCS}}$ \PMcsMlMedTwinPrelx{} under median loss and \PMcsMlMeanTwinPrel{} under winsorized mean loss. Neither model beats persistence.

\paragraph{Inference with few years} The confirmatory results of Section~\ref{sec:pillarA} survive inference that treats years as the only independent units. For S1 the wild-cluster bootstrap by year (Webb weights) gives $p$ \PWildSonePrel{} and the exact randomization test $p$ \PRandSonePrel. The second of these is the smallest value possible with eleven years. For the simulator's encompassing coefficient the bootstrap gives $p$ \PWildEncPrel{}, and the coefficient is positive in \PRandEncPos{} of \PRandEncG{} years (randomization $p$ \PRandEncPrel). For the combination gain the bootstrap gives $p$ \PWildGainPrel{}, and the gain is positive in \PRandGainNeg{} of \PRandGainG{} years (randomization $p$ \PRandGainPrel). Under raw mean loss the model confidence set keeps all \PMcsRawSize{} forecasters, so raw losses cannot separate any of them.

\paragraph{Scale-free errors} Relative to the 16-quarter margin, the simulator's geometric mean relative absolute error is \PGmraeTwin{} (CI \PGmraeTwinLo--\PGmraeTwinHi). The four-quarter margin's is \PGmraeFour{} and the loop-free simulator's \PGmraeNoloops. The simulator beats the margin in \PGmraeTwinBeats{} of firm-years.

Excluding the \PNearZeroN{} firm-years whose trailing revenue was below a quarter of their own 16-quarter level changes the simulator's mean error only from \PCommonMeanTwin{} to \PNearZeroTwinMean{} points (\PCommonMeanN{} firm-years). The explosive firm-years of Table~\ref{tab:tails} are therefore not firms with little revenue relative to their own history. They are firms whose revenue later jumped far above its level at the cutoff.

\paragraph{Stronger density rivals} Two rivals model each firm's own uncertainty. The first scales the margin's residuals by the firm's cash-margin volatility. The second is a quantile regression of the residuals on size, revenue volatility and that volatility. With nominal 90\% intervals they cover \PDenCovHet{} and \PDenCovQr{} of outcomes, against \PDenCovTwin{} for the simulator. Their mean CRPS is \PDenCrpsHet{} and \PDenCrpsQr, against \PDenCrpsTwin{} for the simulator; the heavy tails keep these differences from significance ($p=\PDenHetP$ and \PDenQrP).

The earlier specifications were as narrow: their 90\% coverage is \PDenCovVtwopre{} for v2-pre and \PDenCovVtworev{} for v2-rev. So the corrections did not cause the undercoverage. They did raise the simulator's mean CRPS, through the explosive firm-years.

\paragraph{Recalibrated shortfall probabilities} Isotonic recalibration on earlier cutoffs removes most of the simulator's reliability loss, but not enough. The cleaner label L2 counts petitions, accelerations and the cash criterion, but not delisting notices (\PShortLtwoN{} firm-years, \PShortLtwoEv{} events, nine cutoff years). On it the simulator's Brier skill against the logit rises from \PShortLtwoRaw{} to \PShortLtwoRecal, with a firm-block interval of \PShortLtwoLo{} to \PShortLtwoHi, so the confirmatory test \PShortVerdict. On the original label it reaches \PShortLoneRecal{} (\PShortLoneLo{} to \PShortLoneHi).

In point estimate the recalibrated simulator's resolution (\PShortLtwoRes) equals the logit's (\PShortLtwoResLogit). Once calibrated, its probabilities appear to carry no more information than three ratios.

\paragraph{The NBER look-ahead and the elasticity interval} The plan's opening recession state differs from the one known in real time only at the 2021 cutoff, for \PNberDiffer{} of its \PNberN{} firm-years; at every other cutoff the two states agree. Re-running 2021 with the real-time state leaves the transmission errors unchanged (median \PNberMed{} in both) and the Recession-cell shortfall probability at \PNberShort. Given realized revenue the regime has no effect, and the Recession preset fixes the regime anyway.

The elasticity ratio of v3 has a firm-cluster interval of \PElVthreeFirmLo{} to \PElVthreeFirmHi{} (year-block \PElVthreeYearLo{} to \PElVthreeYearHi). That interval reaches into the pre-specified range, so the failure of the criterion rests on a point estimate, not on an established result. For v2-pre the interval (\PElVtwopreFirmLo--\PElVtwopreFirmHi) lies far above the range.
